% MotiveRec paper draft (v0, 2026-11-13). Venue template not chosen yet: plain article for drafting.
% Numbers: <scratch>/IntuRec/checkpoints/<dataset>/MotiveRec_pilot/pilot_c.json and collab/RESEARCH_LOG_motiverec.md.
% Scope decisions (author): motivation = latent geometry only (no intervention results in this paper);
% novelty = motive grounding (mining, retrieval, integration); the two-level reader is cited as prior design.
\documentclass[10pt,twocolumn]{article}
\usepackage[margin=0.8in]{geometry}
\usepackage{amsmath,amssymb,booktabs,multirow,graphicx,xcolor}
\usepackage[numbers]{natbib}
\usepackage{hyperref}
\newcommand{\todo}[1]{\textcolor{red}{[TODO: #1]}}
\newcommand{\thought}{\mathbf{t}}
\newcommand{\M}{\mathcal{M}}

\title{MotiveRec: Grounding Latent Reasoning in Purchase Motives\\for LLM-based Sequential Recommendation}
\author{\todo{authors}}
\date{}

\begin{document}
\maketitle

\begin{abstract}
Latent reasoning lets an LLM recommender ``think'' in a continuous vector before it names the next item,
avoiding the cost and noise of generated text rationales. We find that, trained only through the
recommendation loss, this latent thought \emph{collapses}: across two Amazon domains the thoughts of
different users have an effective rank of 1--9, the first principal direction carries up to 99.8\% of
their variance, and two random users' thoughts have a cosine similarity of 0.94--1.00. A thought that is
nearly identical for every user cannot carry user-specific reasoning. We argue the cause is that nothing
tells the latent \emph{what} to reason about, and propose \textbf{MotiveRec}, which grounds the latent
thought in the user's \emph{purchase motive}, the reason a person buys. MotiveRec (i) mines an
interpretable vocabulary of motive prototypes from review-grounded purchase rationales, (ii) retrieves
the motives likely to drive a user's next purchase from the interaction history alone, and (iii) builds
the latent thought from the retrieved motives, either from their sentence embeddings or by reading their
text in the LLM's own token space, aligned with the user's mined rationale. Motive grounding removes the
collapse (random-pair cosine 0.54--0.69, effective rank 9--16) \todo{and improves accuracy: main result}.
\end{abstract}

%=====================================================================
\section{Introduction}
LLM-based sequential recommenders cast next-item prediction as generating the title of the next item
from a textual interaction history~\citep{bigrec,tallrec}. Explicit chain-of-thought reasoning before
the prediction helps, but generated rationales are slow and unreliable. \emph{Latent} reasoning replaces
them with one or a few continuous vectors written into the prompt before the answer: LatentR$^3$~\citep{r3}
pools the LLM's hidden states into a thought vector and refines it with reinforcement learning, and
IntuRec~\citep{inturec} builds it from retrieved candidate items (``intuitions'').

These latent thoughts are supervised only by the next-item loss. We ask what they actually encode, and
find that they collapse (Section~\ref{sec:pilot}). For both methods and both datasets, the thoughts of
all test users lie almost on one direction: the effective rank of the thought matrix is 1.0--9.0, the
first principal component explains 31--99.8\% of the variance, and two random users' thoughts have a
cosine of 0.94--1.00. On Toys, R$^3$ produces one constant vector for every user. A latent that is the
same for everyone is a learned prompt, not reasoning about the user.

We attribute the collapse to a missing \emph{content} for the latent: the recommendation loss can be
reduced by a shared, user-independent vector, so nothing pushes the thought to differ between users.
Content-free latents (R$^3$) and item-level content (IntuRec, \emph{what} to buy) both collapse. We
propose to give the latent a different kind of content: the \emph{reason} the user buys. Purchase
motives, such as completing a collection, buying a gift, or loyalty to an artist or franchise, are
interpretable, shared across users and items, and vary between users with the same history items.

\textbf{MotiveRec} grounds latent reasoning in purchase motives in three steps. \emph{Motive mining}
turns review-grounded purchase rationales of training interactions into a compact, named vocabulary of
motive prototypes. \emph{Motive retrieval} predicts, from the history alone, which prototypes are likely
to drive the next purchase, so no target information is used at inference. \emph{Motive integration}
builds the latent thought from the retrieved motives, and an alignment objective ties the thought to the
user's own mined rationale.

\paragraph{Contributions.}
\begin{itemize}
\item We identify and measure \emph{latent collapse} in LLM recommenders with latent reasoning: the
latent thought becomes nearly identical across users (Section~\ref{sec:pilot}).
\item We propose MotiveRec, which grounds the latent in purchase motives through a target-free pipeline
of motive mining, motive retrieval, and motive integration with two variants (embedding and text).
\item We release motive vocabularies for two Amazon domains (859 and 881 named prototypes).
\item Experiments on two datasets show that motive grounding removes the collapse \todo{and improves
recommendation accuracy over R$^3$, IntuRec and the backbone}.
\end{itemize}

%=====================================================================
\section{Pilot study: latent collapse}\label{sec:pilot}
\paragraph{Setup.} We use the Amazon CDs \& Vinyl and Toys \& Games datasets with the same splits
for all models (CDs: 49{,}251 / 6{,}157 train / test rows; Toys: 53{,}898 / 6{,}738), and the original
released code and checkpoints of LatentR$^3$ (after supervised fine-tuning, SFT, and after RL) and
IntuRec, all on Qwen2.5-1.5B. For every test row we take the thought vector $\thought_u\in\mathbb{R}^{d}$
($d=1536$) that the model writes into the prompt.

\paragraph{Measures.} With $T\in\mathbb{R}^{n\times d}$ the matrix of test thoughts and $\sigma_i$ the
singular values of its centred version, we report (i) the \emph{effective rank}
$\exp(-\sum_i p_i\log p_i)$ with $p_i=\sigma_i/\sum_j\sigma_j$~\citep{effrank}, which is $d$ for isotropic
vectors and 1 for vectors on one line; (ii) the share of variance on the first principal component
(PC1); and (iii) the mean cosine similarity of 2{,}000 random pairs of (uncentred) thoughts, which is near
1 when all users receive almost the same vector.

\begin{table}[t]
\centering\small
\caption{Geometry of test-set latent thoughts. Higher effective rank and lower pair cosine mean
more diverse thoughts. \emph{History query}: hidden state of the fine-tuned backbone at the thought position, before the latent is formed (measured on the MotiveRec-T backbone).}
\label{tab:collapse}
\begin{tabular}{llccc}
\toprule
Data & Model & Eff.\ rank & PC1 var. & Pair cos \\
\midrule
\multirow{4}{*}{CDs}
 & LatentR$^3$ (SFT) & 9.01 & 0.307 & 0.944 \\
 & LatentR$^3$ (RL)  & 1.55 & 0.905 & 0.995 \\
 & IntuRec           & 1.02 & 0.998 & 0.952 \\
 & \textit{history query} & \textit{138.0} & \textit{0.140} & \textit{0.686} \\
\midrule
\multirow{4}{*}{Toys}
 & LatentR$^3$ (SFT) & 2.27 & 0.631 & 1.000 \\
 & LatentR$^3$ (RL)  & 2.30 & 0.607 & 1.000 \\
 & IntuRec           & 1.05 & 0.992 & 0.996 \\
 & \textit{history query} & \textit{132.0} & \textit{0.135} & \textit{0.511} \\
\bottomrule
\end{tabular}
\end{table}

\paragraph{Findings.} Table~\ref{tab:collapse} shows that the latent thoughts collapse for every model
and dataset. The pair cosine is 0.94--1.00, and on Toys both R$^3$ checkpoints map every user to the
same vector (the remaining rank is numerical noise). The RL stage of R$^3$ makes the collapse worse on
CDs (effective rank $9.01\to1.55$). IntuRec's thought, built from item intuitions, is almost
one-dimensional on both datasets (PC1 $\geq 0.99$). The collapse is not forced by the backbone: the LLM's hidden state at the thought position is diverse
(effective rank 132--138 in a fine-tuned backbone of the same size), so the user information is available
but is not kept in the latent. \todo{measure the history query of R$^3$ / IntuRec themselves}

\paragraph{Implication.} A latent that is nearly the same for all users cannot express user-specific
reasoning. Since both a content-free latent and a latent built from candidate items collapse, we ask
whether a latent grounded in an explicit, user-varying concept, the purchase motive, stays diverse.

%=====================================================================
\section{MotiveRec}\label{sec:method}
\paragraph{Problem and backbone.} A user $u$ has a chronological history $H_u=(i_1,\dots,i_n)$ of item
titles; the task is to generate the title of the next item $i_{n+1}$. Following latent-reasoning
recommenders~\citep{r3,inturec}, the prompt contains the history followed by a reserved thought token.
A first pass of the LLM gives the hidden state at the thought position, the \emph{history query}
$\mathbf{h}_u\in\mathbb{R}^d$. The thought vector $\thought_u$ replaces the thought token's input
embedding, and the LLM then decodes the next title under a catalogue constraint. MotiveRec changes how
$\thought_u$ is formed: from motives.

\subsection{Motive mining}\label{sec:mining}
\paragraph{Purchase rationales.} For every \emph{training} interaction, an LLM (Qwen3-8B) reads the
history, the purchased item and the user's review of it, and writes 1--2 short motive phrases (2--6
words) and a one-sentence first-person rationale (``I bought this to \dots''). Reviews make the
rationale grounded in what the buyer actually said rather than in item attributes alone.

\paragraph{From phrases to prototypes.} Motive phrases are open-vocabulary and highly redundant
(``60s rock collection'', ``classic 60's music''). We embed the distinct phrases (Qwen3-Embedding with a
purchase-motive instruction), build a similarity graph and cluster it with Louvain community
detection~\citep{louvain} (about 1{,}000 communities per domain). Near-duplicate communities (centroid cosine
above a threshold) are merged, and communities supported by too few training rows are folded into their
nearest neighbour. An LLM then names each community from its most frequent phrases with a canonical
motive phrase, a first-person reason and a label, and flags communities that describe post-purchase
complaints rather than motives; these are removed. This yields a vocabulary
$\M=\{m_1,\dots,m_K\}$ of $K=859$ (CDs) and $K=881$ (Toys) named motive prototypes, each with text
$x_k=$ ``\emph{motive}: \emph{reason}'' (e.g., ``Exploring 60s Music: I want to discover the sounds and
styles of the 1960s music era.'').

\paragraph{Motive labels.} Each phrase of each interaction is assigned to its nearest prototype, so each
training interaction $j$ receives a label set $Y_j\subseteq\M$ of 0--2 motives (99.5\% of rows have at
least one). Labels are used only as training targets of the retriever and are never model inputs.

\subsection{Motive retrieval}\label{sec:retrieval}
At test time the target item, and hence its review and rationale, is unknown. A retriever therefore
predicts the motives of the next purchase from the history only. Its input concatenates
(i) a SASRec~\citep{sasrec} embedding of the item sequence, (ii) a sentence embedding (gte-large) of the
history titles, (iii) an \emph{item motive profile}, the average over history items of the motive
distribution of all training purchases of that item (why people buy these items), and (iv) the motive
distribution of the user's own earlier training interactions. An MLP scores all $K$ prototypes and is
trained with a multi-positive softmax cross-entropy on $Y_j$, with early stopping on validation MRR.
For training rows of the recommender, retrieved motives come from 5-fold cross-fitting grouped by user, so
no row's motives are predicted by a retriever that saw its own label. MotiveRec uses the top-$k$
prototypes $R_u=(m_{r_1},\dots,m_{r_k})$, $k=10$.
% Retriever quality (test): Recall@10 0.42 (CDs) / 0.32 (Toys); MRR 0.32 / 0.29; popularity prior 0.38 / 0.26.

\subsection{Motive integration}\label{sec:integration}
The retrieved motives are turned into the latent thought by a history-conditioned selection: the
user's history decides which candidate motives matter. We study two ways to represent a motive.

\paragraph{Embedding integration (MotiveRec-E).} Each prototype text is encoded once by a frozen
sentence encoder, $\mathbf{e}_k=\mathrm{Enc}(x_k)$, in the same space as the mined rationales. The motives
are projected into the LLM space, $\mathbf{s}_j=W_p\,\mathbf{e}_{r_j}$, and attended by the history
query with a residual path:
\begin{equation}
\thought_u = W_r\mathbf{h}_u + \sum_{j=1}^{k}\alpha_j\,W_v\mathbf{s}_j,\quad
\alpha=\mathrm{softmax}_j\!\left(\frac{(W_q\mathbf{h}_u)^{\top}W_k\mathbf{s}_j}{\sqrt d}\right).
\end{equation}
The residual keeps the history signal, and the attention weights $\alpha$ give an interpretable
distribution over the retrieved motives.

\paragraph{Text integration (MotiveRec-T).} Instead of an external encoder, the LLM reads each motive in
its own vocabulary, so that motives live in the same embedding space as the tokens the decoder consumes.
A motive text $x_{r_j}$ is tokenised and embedded with the LLM's input embedding table,
$X_j\in\mathbb{R}^{L\times d}$. A \emph{motive reader} compresses it in two levels, adopting the two-level
attention pooling design of~\citet{inturec}:
\begin{align}
\mathbf{z}_j &= \mathrm{softmax}\!\left(\frac{\mathbf{q}_{j,L}^{\top}K_j^{\top}}{\sqrt d}\right)V_j,\quad
[Q_j,K_j,V_j]=\mathrm{RoPE}(X_jW^{(1)}_{q,k,v}),\\
\thought_u &= \mathrm{softmax}_j\!\left(\frac{(W^{(2)}_q\mathbf{h}_u)^{\top}W^{(2)}_k\mathbf{z}_j}{\sqrt d}\right)W^{(2)}_v\mathbf{z}_j.
\end{align}
The first level summarises each motive into one vector with its last-token query and rotary positions;
the second level selects among the $k$ motive vectors according to the user's history. The only new
parameters are the reader's projections; motive texts share the LLM's embeddings.

\subsection{Motive-aligned training}\label{sec:training}
The recommendation loss alone allows a collapsed latent (Section~\ref{sec:pilot}). We add an alignment
term that asks the thought to identify the user's own purchase rationale $\rho_u$ (Section~\ref{sec:mining};
available for training rows only) among $N$ random training rationales $\{\rho_n^-\}$:
\begin{equation}
\mathcal{L}_{\text{align}} = -\log\frac{\exp(\langle g(\thought_u),\, \phi(\rho_u)\rangle/\tau)}
{\exp(\langle g(\thought_u),\phi(\rho_u)\rangle/\tau)+\sum_{n}\exp(\langle g(\thought_u),\phi(\rho^-_n)\rangle/\tau)},
\end{equation}
with cosine similarity $\langle\cdot,\cdot\rangle$ and temperature $\tau=0.05$. For MotiveRec-E, $\phi$ is the
frozen sentence encoder and $g$ a linear head; for MotiveRec-T, $\phi$ is the first level of the motive
reader applied to the rationale text and $g$ the identity, so no external embedding space is involved.
The full objective is $\mathcal{L}=\mathcal{L}_{\text{NTP}}+\beta\,\mathcal{L}_{\text{align}}$ with
$\beta=0.1$, where $\mathcal{L}_{\text{NTP}}$ is the next-token loss on the target title. Rationales are
used only as training targets; at inference MotiveRec needs only the history and the retrieved motives.

%=====================================================================
\section{Experiments}\label{sec:exp}
\paragraph{Datasets and protocol.} \todo{dataset statistics table; leave-last-out split; constrained beam
decoding (10 beams); HR@$\{1,5,10\}$ and NDCG@$\{5,10\}$ on the full test set.}

\paragraph{Baselines.} \todo{SASRec, GRU4Rec; BIGRec / D3; backbone with an empty thought slot
(same training); LatentR$^3$ (SFT, RL); IntuRec.}

\paragraph{Implementation.} Qwen2.5-1.5B backbone, learning rate $3\times10^{-4}$, effective batch 128,
$k=10$ motives, $\beta=0.1$, $\tau=0.05$, $N$ negatives \todo{63 (T) / 1024 (E)}.

\subsection{Main results}
\todo{Table: MotiveRec-E / MotiveRec-T vs.\ baselines on both datasets.}

\subsection{Does motive grounding prevent collapse?}
\begin{table}[t]
\centering\small
\caption{Thought geometry with motive grounding (MotiveRec-T) vs.\ Table~\ref{tab:collapse}.}
\label{tab:geometry}
\begin{tabular}{llccc}
\toprule
Data & Model & Eff.\ rank & PC1 var. & Pair cos \\
\midrule
\multirow{3}{*}{CDs}  & LatentR$^3$ (RL) & 1.55 & 0.905 & 0.995 \\
                      & IntuRec          & 1.02 & 0.998 & 0.952 \\
                      & MotiveRec-T      & 8.78 & 0.353 & 0.687 \\
\midrule
\multirow{3}{*}{Toys} & LatentR$^3$ (RL) & 2.30 & 0.607 & 1.000 \\
                      & IntuRec          & 1.05 & 0.992 & 0.996 \\
                      & MotiveRec-T      & 15.53 & 0.200 & 0.540 \\
\bottomrule
\end{tabular}
\end{table}
MotiveRec's thoughts differ between users (pair cosine 0.54--0.69, close to the diversity of the history
query itself) and spread over many directions (PC1 0.20--0.35). They are also organised by
behaviour: on Toys, the centred cosine between a test thought and the training thoughts of the same user
is 0.44 and of training rows sharing a motive is 0.10, against 0.006 for random pairs.
\todo{MotiveRec-E geometry; ablation $\beta=0$ to separate the effect of motive input and alignment.}

\subsection{Ablations and analysis}
\todo{(i) number of motives $k$; (ii) $\beta=0$; (iii) random motives instead of retrieved ones;
(iv) MotiveRec-E vs.\ -T; (v) retriever quality vs.\ accuracy; (vi) case studies with attention weights
over retrieved motives.}

%=====================================================================
\section{Related work}
\paragraph{LLM-based recommendation.} \todo{TALLRec, BIGRec, D3, LC-Rec, \dots}
\paragraph{Reasoning for recommendation.} Explicit rationales \todo{cite}; latent reasoning:
LatentR$^3$~\citep{r3} learns a content-free latent and refines it with RL; IntuRec~\citep{inturec}
conditions the latent on retrieved candidate items, i.e., on \emph{what} the user may buy. MotiveRec
conditions it on \emph{why} the user buys, with an explicit, mined and interpretable motive vocabulary.
\paragraph{Intent- and motive-aware recommendation.} \todo{latent intent models (ICLRec, DSSRec), review-based
user modelling; difference: open-vocabulary, LLM-mined, named motives used to ground LLM reasoning.}
\paragraph{Representation collapse.} \todo{dimensional collapse in contrastive / recommendation
embeddings; effective rank.}

\section{Conclusion}
\todo{}

\bibliographystyle{plainnat}
\bibliography{refs}
\end{document}
